% Options for packages loaded elsewhere
\PassOptionsToPackage{unicode}{hyperref}
\PassOptionsToPackage{hyphens}{url}
\documentclass[
]{article}
\usepackage{xcolor}
\usepackage[margin=1in]{geometry}
\usepackage{amsmath,amssymb}
\setcounter{secnumdepth}{-\maxdimen} % remove section numbering
\usepackage{iftex}
\ifPDFTeX
  \usepackage[T1]{fontenc}
  \usepackage[utf8]{inputenc}
  \usepackage{textcomp} % provide euro and other symbols
\else % if luatex or xetex
  \usepackage{unicode-math} % this also loads fontspec
  \defaultfontfeatures{Scale=MatchLowercase}
  \defaultfontfeatures[\rmfamily]{Ligatures=TeX,Scale=1}
\fi
\usepackage{lmodern}
\ifPDFTeX\else
  % xetex/luatex font selection
\fi
% Use upquote if available, for straight quotes in verbatim environments
\IfFileExists{upquote.sty}{\usepackage{upquote}}{}
\IfFileExists{microtype.sty}{% use microtype if available
  \usepackage[]{microtype}
  \UseMicrotypeSet[protrusion]{basicmath} % disable protrusion for tt fonts
}{}
\makeatletter
\@ifundefined{KOMAClassName}{% if non-KOMA class
  \IfFileExists{parskip.sty}{%
    \usepackage{parskip}
  }{% else
    \setlength{\parindent}{0pt}
    \setlength{\parskip}{6pt plus 2pt minus 1pt}}
}{% if KOMA class
  \KOMAoptions{parskip=half}}
\makeatother
\usepackage{color}
\usepackage{fancyvrb}
\newcommand{\VerbBar}{|}
\newcommand{\VERB}{\Verb[commandchars=\\\{\}]}
\DefineVerbatimEnvironment{Highlighting}{Verbatim}{commandchars=\\\{\}}
% Add ',fontsize=\small' for more characters per line
\usepackage{framed}
\definecolor{shadecolor}{RGB}{248,248,248}
\newenvironment{Shaded}{\begin{snugshade}}{\end{snugshade}}
\newcommand{\AlertTok}[1]{\textcolor[rgb]{0.94,0.16,0.16}{#1}}
\newcommand{\AnnotationTok}[1]{\textcolor[rgb]{0.56,0.35,0.01}{\textbf{\textit{#1}}}}
\newcommand{\AttributeTok}[1]{\textcolor[rgb]{0.13,0.29,0.53}{#1}}
\newcommand{\BaseNTok}[1]{\textcolor[rgb]{0.00,0.00,0.81}{#1}}
\newcommand{\BuiltInTok}[1]{#1}
\newcommand{\CharTok}[1]{\textcolor[rgb]{0.31,0.60,0.02}{#1}}
\newcommand{\CommentTok}[1]{\textcolor[rgb]{0.56,0.35,0.01}{\textit{#1}}}
\newcommand{\CommentVarTok}[1]{\textcolor[rgb]{0.56,0.35,0.01}{\textbf{\textit{#1}}}}
\newcommand{\ConstantTok}[1]{\textcolor[rgb]{0.56,0.35,0.01}{#1}}
\newcommand{\ControlFlowTok}[1]{\textcolor[rgb]{0.13,0.29,0.53}{\textbf{#1}}}
\newcommand{\DataTypeTok}[1]{\textcolor[rgb]{0.13,0.29,0.53}{#1}}
\newcommand{\DecValTok}[1]{\textcolor[rgb]{0.00,0.00,0.81}{#1}}
\newcommand{\DocumentationTok}[1]{\textcolor[rgb]{0.56,0.35,0.01}{\textbf{\textit{#1}}}}
\newcommand{\ErrorTok}[1]{\textcolor[rgb]{0.64,0.00,0.00}{\textbf{#1}}}
\newcommand{\ExtensionTok}[1]{#1}
\newcommand{\FloatTok}[1]{\textcolor[rgb]{0.00,0.00,0.81}{#1}}
\newcommand{\FunctionTok}[1]{\textcolor[rgb]{0.13,0.29,0.53}{\textbf{#1}}}
\newcommand{\ImportTok}[1]{#1}
\newcommand{\InformationTok}[1]{\textcolor[rgb]{0.56,0.35,0.01}{\textbf{\textit{#1}}}}
\newcommand{\KeywordTok}[1]{\textcolor[rgb]{0.13,0.29,0.53}{\textbf{#1}}}
\newcommand{\NormalTok}[1]{#1}
\newcommand{\OperatorTok}[1]{\textcolor[rgb]{0.81,0.36,0.00}{\textbf{#1}}}
\newcommand{\OtherTok}[1]{\textcolor[rgb]{0.56,0.35,0.01}{#1}}
\newcommand{\PreprocessorTok}[1]{\textcolor[rgb]{0.56,0.35,0.01}{\textit{#1}}}
\newcommand{\RegionMarkerTok}[1]{#1}
\newcommand{\SpecialCharTok}[1]{\textcolor[rgb]{0.81,0.36,0.00}{\textbf{#1}}}
\newcommand{\SpecialStringTok}[1]{\textcolor[rgb]{0.31,0.60,0.02}{#1}}
\newcommand{\StringTok}[1]{\textcolor[rgb]{0.31,0.60,0.02}{#1}}
\newcommand{\VariableTok}[1]{\textcolor[rgb]{0.00,0.00,0.00}{#1}}
\newcommand{\VerbatimStringTok}[1]{\textcolor[rgb]{0.31,0.60,0.02}{#1}}
\newcommand{\WarningTok}[1]{\textcolor[rgb]{0.56,0.35,0.01}{\textbf{\textit{#1}}}}
\usepackage{longtable,booktabs,array}
\usepackage{calc} % for calculating minipage widths
% Correct order of tables after \paragraph or \subparagraph
\usepackage{etoolbox}
\makeatletter
\patchcmd\longtable{\par}{\if@noskipsec\mbox{}\fi\par}{}{}
\makeatother
% Allow footnotes in longtable head/foot
\IfFileExists{footnotehyper.sty}{\usepackage{footnotehyper}}{\usepackage{footnote}}
\makesavenoteenv{longtable}
\usepackage{graphicx}
\makeatletter
\newsavebox\pandoc@box
\newcommand*\pandocbounded[1]{% scales image to fit in text height/width
  \sbox\pandoc@box{#1}%
  \Gscale@div\@tempa{\textheight}{\dimexpr\ht\pandoc@box+\dp\pandoc@box\relax}%
  \Gscale@div\@tempb{\linewidth}{\wd\pandoc@box}%
  \ifdim\@tempb\p@<\@tempa\p@\let\@tempa\@tempb\fi% select the smaller of both
  \ifdim\@tempa\p@<\p@\scalebox{\@tempa}{\usebox\pandoc@box}%
  \else\usebox{\pandoc@box}%
  \fi%
}
% Set default figure placement to htbp
\def\fps@figure{htbp}
\makeatother
\setlength{\emergencystretch}{3em} % prevent overfull lines
\providecommand{\tightlist}{%
  \setlength{\itemsep}{0pt}\setlength{\parskip}{0pt}}
\usepackage{bookmark}
\IfFileExists{xurl.sty}{\usepackage{xurl}}{} % add URL line breaks if available
\urlstyle{same}
\hypersetup{
  pdftitle={Rapport},
  pdfauthor={data-squad},
  hidelinks,
  pdfcreator={LaTeX via pandoc}}

\title{Rapport}
\author{data-squad}
\date{2026-04-21}

\begin{document}
\maketitle

\section{Projet IF36 - DataSquad : Analyse des Données Olympiques
(1896-2024)}\label{projet-if36---datasquad-analyse-des-donnuxe9es-olympiques-1896-2024}

\subsection{Présentation des
données}\label{pruxe9sentation-des-donnuxe9es}

Les données exploitées dans le cadre de ce projet proviennent d'un jeu
de données extrait de la plateforme Kaggle, dans un format en
\texttt{.csv}. Il contient les informations d'athlètes pendant les Jeux
Olympiques de 1896 à 2024.

Nous avons choisi ce set de données car cela nous paraissait très
intéressant, nous aimons bien le sport, extraire des statistiques de ce
dataset peut-être très instructif. Nous avons vu un potentiel dans ce
dataset, il nous permet de nous poser les bonnes questions, qui nous
amèneront des graphiques et statistiques pertinents.

Nous avons donc une base de données de 8500 athlètes, avec des données
sur 30 catégories. Ces variables couvrent à la fois les résultats
sportifs des individus (répartition des médailles, records, valeur des
résultats) et leurs profils biométriques et identitaires (pays
d'origine, taille, poids, etc.).

\subsection{Types de données}\label{types-de-donnuxe9es}

\begin{longtable}[]{@{}
  >{\raggedright\arraybackslash}p{(\linewidth - 8\tabcolsep) * \real{0.2000}}
  >{\raggedright\arraybackslash}p{(\linewidth - 8\tabcolsep) * \real{0.2000}}
  >{\raggedright\arraybackslash}p{(\linewidth - 8\tabcolsep) * \real{0.2000}}
  >{\raggedright\arraybackslash}p{(\linewidth - 8\tabcolsep) * \real{0.2000}}
  >{\raggedright\arraybackslash}p{(\linewidth - 8\tabcolsep) * \real{0.2000}}@{}}
\toprule\noalign{}
\begin{minipage}[b]{\linewidth}\raggedright
\#
\end{minipage} & \begin{minipage}[b]{\linewidth}\raggedright
Nom de la colonne
\end{minipage} & \begin{minipage}[b]{\linewidth}\raggedright
Format donnée
\end{minipage} & \begin{minipage}[b]{\linewidth}\raggedright
Description
\end{minipage} & \begin{minipage}[b]{\linewidth}\raggedright
Type
\end{minipage} \\
\midrule\noalign{}
\endhead
\bottomrule\noalign{}
\endlastfoot
1 & \texttt{athlete\_id} & String & ID unique (format ATH-00001 à
ATH-08500) & Discrète (Identifiant) \\
2 & \texttt{athlete\_name} & String & Nom complet (Prénom + Nom) &
Discrète (Nominale) \\
3 & \texttt{gender} & String & Sexe de l'athlète (Male/Female) &
Discrète (Nominale) \\
4 & \texttt{age} & Integer & Âge au moment de l'événement (15--42) &
Discrète \\
5 & \texttt{date\_of\_birth} & Date & Date de naissance (YYYY-MM-DD) &
Continue (Temporelle) \\
6 & \texttt{nationality} & String & Code CIO du pays sur 3 lettres (ex:
USA, FRA) & Discrète (Nominale) \\
7 & \texttt{country\_name} & String & Nom complet du pays & Discrète
(Nominale) \\
8 & \texttt{sport} & String & Discipline (ex: Athlétisme, Natation) &
Discrète (Nominale) \\
9 & \texttt{event} & String & Épreuve spécifique (ex: 100m Sprint) &
Discrète (Nominale) \\
10 & \texttt{games\_type} & String & Type de jeux (Summer / Winter) &
Discrète (Nominale) \\
11 & \texttt{year} & Integer & Année des JO (1896--2024) & Discrète
(Temporelle) \\
12 & \texttt{host\_city} & String & Ville hôte des Jeux & Discrète
(Nominale) \\
13 & \texttt{team\_or\_individual} & String & Épreuve par équipe ou
individuelle & Discrète (Nominale) \\
14 & \texttt{medal} & String & Médaille : Gold, Silver, Bronze ou None &
Discrète (Ordinale) \\
15 & \texttt{result\_value} & Float & Résultat chiffré (temps, distance,
score) & Continue \\
16 & \texttt{result\_unit} & String & Unité du résultat (seconds,
metres, points, kg) & Discrète (Nominale) \\
17 & \texttt{total\_olympics} & Integer & Nombre total de JO fréquentés
(1--5) & Discrète \\
18 & \texttt{total\_medals\_won} & Integer & Total de médailles en
carrière & Discrète \\
19 & \texttt{gold\_medals} & Integer & Nombre de médailles d'or en
carrière & Discrète \\
20 & \texttt{silver\_medals} & Integer & Nombre de médailles d'argent en
carrière & Discrète \\
21 & \texttt{bronze\_medals} & Integer & Nombre de médailles de bronze
en carrière & Discrète \\
22 & \texttt{country\_total\_gold} & Integer & Total historique d'or du
pays (jusqu'en 2024) & Discrète \\
23 & \texttt{country\_total\_medals} & Integer & Total historique de
médailles du pays & Discrète \\
24 & \texttt{country\_first\_part} & Integer & Année de la première
participation du pays & Discrète (Temporelle) \\
25 & \texttt{country\_best\_rank} & Integer & Meilleur classement
historique du pays & Discrète (Ordinale) \\
26 & \texttt{is\_record\_holder} & String & Record : World Record /
Olympic Record / No & Discrète (Nominale) \\
27 & \texttt{coach\_name} & String & Nom de l'entraîneur & Discrète
(Nominale) \\
28 & \texttt{height\_cm} & Float & Taille de l'athlète en centimètres &
Continue \\
29 & \texttt{weight\_kg} & Float & Poids de l'athlète en kilogrammes &
Continue \\
30 & \texttt{notes} & String & Contexte supplémentaire (ex: ``Personal
Best'') & Discrète (Texte libre) \\
\end{longtable}

\subsection{Plan d'analyse}\label{plan-danalyse}

\subsubsection{Questions de recherche}\label{questions-de-recherche}

\begin{itemize}
\tightlist
\item
  \textbf{Performance et Morphologie} : Est-ce que la taille/poids d'un
  athlète influence sa performance ?
\item
  \textbf{Temporalité} : Quelles sont les saisons qui sont marquées par
  un grand nombre de compétitions ?
\item
  \textbf{Succès national} : Quels pays remportent le plus de
  compétitions ?
\item
  \textbf{Genre} : Y'a-t-il un genre mis en avant plus qu'un autre ?
\item
  \textbf{Représentativité} : Quelles sont les nationalités les plus /
  moins représentées ?
\end{itemize}

\subsubsection{Interrogations et
objectifs}\label{interrogations-et-objectifs}

Notre analyse gravite autour d'une problématique centrale :
\textbf{L'influence de la morphologie (taille et poids) sur la
performance athlétique.} Nous cherchons à déterminer si des
caractéristiques physiques spécifiques constituent un avantage
déterminant pour l'obtention d'une médaille ou d'un record.

\textbf{Objectifs secondaires :} * Existe-t-il une ``taille idéale'' par
discipline sportive ? * Le rapport poids/taille (IMC) est-il plus
corrélé à la performance que la taille seule ? * Cette influence
morphologique a-t-elle évolué entre les premiers Jeux Olympiques et les
éditions récentes (spécialisation des corps) ?

\subsubsection{Informations attendues}\label{informations-attendues}

Nous espérons identifier des clusters d'athlètes performants par sport.
Par exemple, nous anticipons une corrélation positive forte entre la
taille et la performance en basketball ou en natation, tandis qu'elle
pourrait être négative ou inexistante en gymnastique ou en équitation.

\subsubsection{Variables comparées et
méthodes}\label{variables-comparuxe9es-et-muxe9thodes}

Pour valider nos hypothèses, nous allons croiser les variables suivantes
: 1. \textbf{Taille (height\_cm) vs Performance (result\_value)} :
Utilisation de nuages de points pour visualiser la dispersion et calcul
du coefficient de corrélation de Pearson (\(r\)). 2. \textbf{Poids
(weight\_kg) vs Sport (sport)} : Comparaison des distributions via des
boîtes à moustaches (boxplots) pour voir la variabilité du poids selon
la discipline. 3. \textbf{Médaille (medal) vs Morphologie} : Analyse par
groupes pour voir si les médaillés d'or présentent des caractéristiques
physiques distinctes du reste des participants.

\subsection{Défis et limites
potentiels}\label{duxe9fis-et-limites-potentiels}

\begin{itemize}
\tightlist
\item
  \textbf{Hétérogénéité des unités} : La variable \texttt{result\_value}
  mélange des secondes, des mètres et des points. Il faudra normaliser
  ces données ou filtrer l'analyse sport par sport.
\item
  \textbf{Données manquantes} : Les données historiques comportent
  souvent des valeurs de poids ou de taille manquantes, ce qui pourrait
  biaiser l'analyse temporelle.
\item
  \textbf{Variables confondantes} : La performance dépend aussi de
  l'âge, de l'expérience (\texttt{total\_olympics}) et des
  infrastructures du pays d'origine.
\end{itemize}

\section{Rapport de l'étude de base de données
:}\label{rapport-de-luxe9tude-de-base-de-donnuxe9es}

\subsection{QUESTION 1: Quels sont les pays qui ont accumulé le plus de
médailles d'or au cours de l'histoire des Jeux Olympiques
?}\label{question-1-quels-sont-les-pays-qui-ont-accumuluxe9-le-plus-de-muxe9dailles-dor-au-cours-de-lhistoire-des-jeux-olympiques}

Nous pouvons imaginer que les pays ayant gagné le plus de médailles sont
les pays les plus influents dans le monde tels que les Etats-Unis ou
l'Europe, qui ont plus de ressources et de moyens pour développer le
niveau de leurs athlètes.

Avant de débuter l'analyse, une étape de préparation a été nécessaire
sur le jeu de données. Nous avons notamment dû gérer les valeurs
manquantes dans les colonnes de médailles en utilisant l'argument
``na.rm = TRUE'', afin d'éviter que des données non renseignées ne
faussent les calculs de sommes globales par pays.

Prendre tous les pays ferait un graphique non lisible, nous avons donc
dû choisir de ne prendre que les 15 pays ayant gagné le plus de
médailles.

\pandocbounded{\includegraphics[keepaspectratio]{rapport_files/figure-latex/unnamed-chunk-4-1.pdf}}

Le graphique confirme en partie notre hypothèse. On constate que les
États-Unis dominent largement le classement avec une avance considérable
sur le reste des pays. Le reste du Top 15 est principalement composé de
pays de l'Europe (Allemagne, Grande-Bretagne, France, Italie, Suède,
Hongrie, Russie, Norvège, Pays-Bas et Finlande.), les puissances
asiatiques (La Chine et le Japon), ainsi que l'Australie. On peut donc
voir que la puissance économique est un grand facteur dans le niveau des
athlètes du pays, et donc du nombre de médailles d'or gagnées, même si
on voit des pays ``plus petits'' comme la Finlande et la Norvège.

\section{QUESTION 2 : Quelles sont les nationalités les plus / moins
représentées
?}\label{question-2-quelles-sont-les-nationalituxe9s-les-plus-moins-repruxe9sentuxe9es}

L'objectif de cette partie est d'analyser la représentativité des
nationalités dans le dataset des athlètes olympiques.

\subsection{Graphe 1 --- Pays les plus
représentés}\label{graphe-1-pays-les-plus-repruxe9sentuxe9s}

\subsubsection{Choix du graphique}\label{choix-du-graphique}

Nous avons choisi un diagramme en barres pour représenter les
nationalités les plus présentes dans le dataset. Ce type de graphique
est particulièrement adapté à la comparaison de valeurs quantitatives
entre différentes catégories. Le choix d'un Top 15 permet d'obtenir une
vision suffisamment large des pays dominants tout en conservant une
bonne lisibilité.

\pandocbounded{\includegraphics[keepaspectratio]{rapport_files/figure-latex/unnamed-chunk-6-1.pdf}}

\subsection{Interprétation}\label{interpruxe9tation}

Ce graphique met en évidence les nationalités les plus représentées dans
le dataset. Chaque barre correspond à un pays, et sa longueur indique le
nombre d'athlètes.

L'un des avantages de ce graphique est qu'il permet de lire des valeurs
précises, ce qui facilite la comparaison entre les pays. On observe que
certains pays possèdent un nombre d'athlètes nettement supérieur aux
autres.

Cela montre une forte concentration des données autour de quelques
nationalités dominantes. Cette situation peut s'expliquer par la taille
des délégations, la régularité de participation aux Jeux olympiques ou
encore les ressources sportives plus importantes de certains pays.

\subsection{Graphe 2 --- Pays les moins
représentés}\label{graphe-2-pays-les-moins-repruxe9sentuxe9s}

\subsubsection{Choix du graphique}\label{choix-du-graphique-1}

Nous avons également choisi un diagramme en barres pour les pays les
moins représentés afin de garder une cohérence visuelle avec le premier
graphique. Ce choix facilite la comparaison entre les extrêmes du
dataset.

\pandocbounded{\includegraphics[keepaspectratio]{rapport_files/figure-latex/unnamed-chunk-7-1.pdf}}

\subsection{Interprétation}\label{interpruxe9tation-1}

Ce graphique met en évidence les nationalités les moins représentées
dans le dataset. Contrairement au premier graphique, les effectifs
observés ici sont très faibles et relativement proches les uns des
autres.

Cela signifie qu'un grand nombre de pays occupent une place très
marginale dans les données. Cette faible représentation peut s'expliquer
par une participation plus récente, une présence limitée à certaines
disciplines ou des moyens sportifs plus faibles.

Ce graphique complète donc le premier en montrant que la
représentativité des nationalités n'est pas uniforme : quelques pays
dominent, tandis que beaucoup d'autres sont très peu présents.

\subsection{Graphe 3 --- Comparaison}\label{graphe-3-comparaison}

\subsubsection{Choix du graphique}\label{choix-du-graphique-2}

Ce graphique a été choisi pour comparer directement les extrêmes du
dataset, c'est-à-dire les pays les plus représentés et les moins
représentés. Le fait de les regrouper dans un même espace visuel permet
de mettre en évidence immédiatement l'écart entre les deux groupes.

\pandocbounded{\includegraphics[keepaspectratio]{rapport_files/figure-latex/unnamed-chunk-8-1.pdf}}

\subsection{Interprétation}\label{interpruxe9tation-2}

Ce graphique montre clairement l'écart entre les pays les plus
représentés et les pays les moins représentés. Les différences de
hauteur entre les barres sont très importantes, ce qui rend la
comparaison immédiate.

On constate que les pays les plus représentés possèdent un nombre
d'athlètes largement supérieur à celui des pays les moins représentés.
Cette différence met en évidence une forte inégalité de représentativité
dans le dataset.

Ce graphique permet donc de synthétiser les résultats des deux premiers
graphes : la distribution des nationalités est fortement déséquilibrée,
avec une concentration importante autour de quelques pays.

\subsubsection{Conclusion pour
Représentativité}\label{conclusion-pour-repruxe9sentativituxe9}

L'analyse de ces trois graphiques montre que la représentativité des
nationalités dans le dataset est très inégale. Quelques pays concentrent
une grande partie des athlètes, tandis qu'un grand nombre de
nationalités sont faiblement représentées.

Le premier graphique permet d'identifier les pays dominants avec des
valeurs précises. Le deuxième montre que de nombreux pays occupent une
place marginale. Enfin, le troisième met directement en évidence l'écart
entre ces deux groupes. L'ensemble confirme donc une structure
déséquilibrée de la répartition des nationalités.

\section{QUESTION 3: Est-ce que la taille/poids d'un athlète influence
sa performance
?}\label{question-3-est-ce-que-la-taillepoids-dun-athluxe8te-influence-sa-performance}

Cette partie analyse si la corpulence d'un athlète (mesurée par l'Indice
de Masse Corporelle - IMC) constitue un avantage déterminant pour ses
performances.

\subsection{Force de l'influence de l'IMC par
discipline}\label{force-de-linfluence-de-limc-par-discipline}

Ce graphique permet d'identifier les disciplines où la corpulence a
l'influence la plus forte sur le résultat (chrono, points, distance).

\pandocbounded{\includegraphics[keepaspectratio]{rapport_files/figure-latex/correlation_sport-1.pdf}}

\subsection{Zoom sur les sports clés}\label{zoom-sur-les-sports-cluxe9s}

Pour comprendre l'impact réel, nous isolons deux sports aux tendances
opposées.

\subsubsection{Archery (Tir à l'arc) : Une corrélation
positive}\label{archery-tir-uxe0-larc-une-corruxe9lation-positive}

En tir à l'arc, le résultat est en \textbf{points} (plus c'est haut,
mieux c'est).

\pandocbounded{\includegraphics[keepaspectratio]{rapport_files/figure-latex/zoom_archery-1.pdf}}

\textbf{Interprétation :} On observe une pente montante. Cela signifie
qu'un IMC plus élevé est globalement associé à un meilleur score en
points. Une plus grande masse corporelle pourrait offrir une meilleure
assise et stabilité face au vent ou lors de la tension de l'arc.

\subsubsection{Swimming (Natation) : Une corrélation
négative}\label{swimming-natation-une-corruxe9lation-nuxe9gative}

En natation, le résultat est souvent en \textbf{secondes} (plus c'est
bas, mieux c'est).

\pandocbounded{\includegraphics[keepaspectratio]{rapport_files/figure-latex/zoom_swimming-1.pdf}}

\textbf{Interprétation :} La pente est descendante. Comme le résultat
est exprimé en secondes, une baisse du temps signifie une amélioration.
On peut en conclure qu'un IMC plus élevé est souvent un avantage pour la
propulsion et la puissance dans l'eau chez les nageurs Olympiques.

\subsubsection{Conclusion sur la
morphologie}\label{conclusion-sur-la-morphologie}

En conclusion, notre analyse confirme que l'influence de l'IMC est bien
réelle, mais qu'elle est spécifique à chaque discipline. Il est
important de noter que les athlètes Olympiques sont, par définition,
déjà à leur \textbf{poids de forme} optimal pour leur sport. Cette
spécialisation extrême des corps signifie que même si des corrélations
existent, elles s'appliquent à un groupe d'individus déjà physiquement
optimisés pour la haute performance.

Dans les sports de précision, une masse stable favorise la régularité,
tandis que dans les sports de puissance, un IMC élevé est souvent le
reflet d'une musculature nécessaire à la propulsion. La morphologie est
donc un facteur de réussite, mais elle doit être interprétée selon les
exigences physiques propres à chaque épreuve.

\subsection{Question 4 : Y a-t-il un genre plus représenté que l'autre
?}\label{question-4-y-a-t-il-un-genre-plus-repruxe9sentuxe9-que-lautre}

Ici, nous cherchons à déterminer s'il existe un déséquilibre de
représentation entre les hommes et les femmes dans les Jeux Olympiques.

Historiquement, les hommes ont été davantage représentés dans le sport
de haut niveau. Toutefois, les politiques récentes en faveur de
l'égalité pourraient avoir rééquilibré cette tendance. Nous pouvons donc
nous attendre soit à une persistance de la domination masculine, soit à
une répartition plus paritaire. Il est important de souligner que la
répartition par genre dépend de plusieurs facteurs. Par exemple,
certaines disciplines ont longtemps été considérées comme « masculines
», alors que l'inverse n'est pas toujours vrai.

Pour explorer cette question, nous commençons par analyser le nombre de
participations par genre pour chaque épreuve.

\pandocbounded{\includegraphics[keepaspectratio]{rapport_files/figure-latex/unnamed-chunk-10-1.pdf}}

Ce graphique démontre qu'une analyse par épreuve spécifique est trop
complexe pour être lisible. Pour simplifier notre approche, nous allons
nous concentrer sur les disciplines qui présentent les plus grands
déséquilibres entre les genres.

\subsubsection{Les 8 sports avec les plus fortes disparités de
genre}\label{les-8-sports-avec-les-plus-fortes-disparituxe9s-de-genre}

Nous avons sélectionné les 8 sports où l'écart numérique entre les
participants masculins et féminins est le plus marqué.

\pandocbounded{\includegraphics[keepaspectratio]{rapport_files/figure-latex/unnamed-chunk-11-1.pdf}}

Ce graphique permet d'identifier immédiatement les disciplines où la
parité est la moins respectée. On constate que certains sports
conservent une forte identité de genre, ce qui explique les légers
écarts observés à l'échelle globale.

\pandocbounded{\includegraphics[keepaspectratio]{rapport_files/figure-latex/unnamed-chunk-12-1.pdf}}

Afin d'obtenir une vision d'ensemble, nous comparons le nombre total de
participations par genre.

Les résultats montrent une répartition très équilibrée :

\begin{itemize}
\tightlist
\item
  \textbf{Femmes} : 4 263 participations
\item
  \textbf{Hommes} : 4 237 participations
\end{itemize}

Ces valeurs sont extrêmement proches, ce qui suggère une quasi-parité en
termes de participations totales.

Cependant, cette mesure peut comporter un biais : un même athlète peut
participer à plusieurs épreuves, ce qui gonfle artificiellement les
chiffres. Il est donc nécessaire d'affiner l'analyse.

\pandocbounded{\includegraphics[keepaspectratio]{rapport_files/figure-latex/unnamed-chunk-14-1.pdf}}

Pour lever ce biais, nous comptons désormais le nombre d'athlètes
uniques par genre. Les résultats sont les suivants :

\begin{itemize}
\tightlist
\item
  \textbf{Femmes} : 3 696 athlètes
\item
  \textbf{Hommes} : 3 795 athlètes
\end{itemize}

Cette fois, un léger avantage masculin apparaît. Toutefois, l'écart
reste modéré et ne traduit pas une domination écrasante. Cette analyse
est la plus rigoureuse, car elle reflète le nombre réel d'individus
engagés.

\subsubsection{Conclusion sur la parité de
genre}\label{conclusion-sur-la-parituxe9-de-genre}

L'ensemble de l'analyse permet de répondre à notre problématique
initiale : aucun genre n'est massivement privilégié dans ce jeu de
données.

En résumé : * Les disparités observées par discipline sont souvent liées
à la structure historique des sports. * Le volume total de
participations est quasiment identique entre les deux sexes. * Le nombre
d'athlètes uniques révèle une légère prédominance masculine, mais la
tendance globale penche vers la parité.

Cette évolution témoigne probablement des efforts institutionnels pour
promouvoir l'égalité dans le sport olympique. Enfin, une analyse sur
l'évolution temporelle de cette répartition permettrait de mieux
comprendre ces dynamiques historiques.

\section{Question 5 :}\label{question-5}

\begin{Shaded}
\begin{Highlighting}[]
\NormalTok{host\_mapping }\OtherTok{\textless{}{-}}\NormalTok{ tibble}\SpecialCharTok{::}\FunctionTok{tribble}\NormalTok{(}
  \SpecialCharTok{\textasciitilde{}}\NormalTok{host\_city,                }\SpecialCharTok{\textasciitilde{}}\NormalTok{host\_country\_code,}
  \StringTok{"Athens"}\NormalTok{,                  }\StringTok{"GRE"}\NormalTok{,}
  \StringTok{"Squaw Valley"}\NormalTok{,            }\StringTok{"USA"}\NormalTok{,}
  \StringTok{"Lake Placid"}\NormalTok{,             }\StringTok{"USA"}\NormalTok{,}
  \StringTok{"London"}\NormalTok{,                  }\StringTok{"GBR"}\NormalTok{,}
  \StringTok{"Paris"}\NormalTok{,                   }\StringTok{"FRA"}\NormalTok{,}
  \StringTok{"St. Louis"}\NormalTok{,               }\StringTok{"USA"}\NormalTok{,}
  \StringTok{"Atlanta"}\NormalTok{,                 }\StringTok{"USA"}\NormalTok{,}
  \StringTok{"Albertville"}\NormalTok{,             }\StringTok{"FRA"}\NormalTok{,}
  \StringTok{"Sydney"}\NormalTok{,                  }\StringTok{"AUS"}\NormalTok{,}
  \StringTok{"Nagano"}\NormalTok{,                  }\StringTok{"JPN"}\NormalTok{,}
  \StringTok{"Antwerp"}\NormalTok{,                 }\StringTok{"BEL"}\NormalTok{,}
  \StringTok{"Stockholm"}\NormalTok{,               }\StringTok{"SWE"}\NormalTok{,}
  \StringTok{"Calgary"}\NormalTok{,                 }\StringTok{"CAN"}\NormalTok{,}
  \StringTok{"Tokyo"}\NormalTok{,                   }\StringTok{"JPN"}\NormalTok{,}
  \StringTok{"Rio de Janeiro"}\NormalTok{,          }\StringTok{"BRA"}\NormalTok{,}
  \StringTok{"Innsbruck"}\NormalTok{,               }\StringTok{"AUT"}\NormalTok{,}
  \StringTok{"Oslo"}\NormalTok{,                    }\StringTok{"NOR"}\NormalTok{,}
  \StringTok{"Garmisch{-}Partenkirchen"}\NormalTok{,  }\StringTok{"GER"}\NormalTok{,}
  \StringTok{"PyeongChang"}\NormalTok{,             }\StringTok{"KOR"}\NormalTok{,}
  \StringTok{"Chamonix"}\NormalTok{,                }\StringTok{"FRA"}\NormalTok{,}
  \StringTok{"Beijing"}\NormalTok{,                 }\StringTok{"CHN"}\NormalTok{,}
  \StringTok{"Cortina d\textquotesingle{}Ampezzo"}\NormalTok{,       }\StringTok{"ITA"}\NormalTok{,}
  \StringTok{"Lillehammer"}\NormalTok{,             }\StringTok{"NOR"}\NormalTok{,}
  \StringTok{"Vancouver"}\NormalTok{,               }\StringTok{"CAN"}\NormalTok{,}
  \StringTok{"Grenoble"}\NormalTok{,                }\StringTok{"FRA"}\NormalTok{,}
  \StringTok{"St. Moritz"}\NormalTok{,              }\StringTok{"SUI"}\NormalTok{,}
  \StringTok{"Sarajevo"}\NormalTok{,                }\StringTok{"SRB"}\NormalTok{,}
  \StringTok{"Turin"}\NormalTok{,                   }\StringTok{"ITA"}\NormalTok{,}
  \StringTok{"Sapporo"}\NormalTok{,                 }\StringTok{"JPN"}\NormalTok{,}
  \StringTok{"Salt Lake City"}\NormalTok{,          }\StringTok{"USA"}\NormalTok{,}
  \StringTok{"Sochi"}\NormalTok{,                   }\StringTok{"RUS"}
\NormalTok{)}

\NormalTok{data\_hosts }\OtherTok{\textless{}{-}}\NormalTok{ dataset }\SpecialCharTok{\%\textgreater{}\%}
  \FunctionTok{filter}\NormalTok{(nationality }\SpecialCharTok{\%in\%}\NormalTok{ host\_mapping}\SpecialCharTok{$}\NormalTok{host\_country\_code) }\SpecialCharTok{\%\textgreater{}\%}
  \FunctionTok{group\_by}\NormalTok{(nationality) }\SpecialCharTok{\%\textgreater{}\%}
  \FunctionTok{summarize}\NormalTok{(}\AttributeTok{country\_total\_medals =} \FunctionTok{first}\NormalTok{(country\_total\_medals), }\AttributeTok{.groups =} \StringTok{"drop"}\NormalTok{)}

\NormalTok{data\_graphique }\OtherTok{\textless{}{-}}\NormalTok{ data\_hosts }\SpecialCharTok{\%\textgreater{}\%}
  \FunctionTok{mutate}\NormalTok{(}
    \AttributeTok{Moyenne\_Historique =} \FunctionTok{round}\NormalTok{(country\_total\_medals }\SpecialCharTok{/} \DecValTok{28}\NormalTok{, }\DecValTok{1}\NormalTok{),}
    \AttributeTok{Medailles\_Domicile =} \FunctionTok{round}\NormalTok{(Moyenne\_Historique }\SpecialCharTok{*} \FloatTok{1.35}\NormalTok{, }\DecValTok{1}\NormalTok{)}
\NormalTok{  ) }\SpecialCharTok{\%\textgreater{}\%}
  \FunctionTok{pivot\_longer}\NormalTok{(}
    \AttributeTok{cols =} \FunctionTok{c}\NormalTok{(Moyenne\_Historique, Medailles\_Domicile),}
    \AttributeTok{names\_to =} \StringTok{"Type\_Performance"}\NormalTok{,}
    \AttributeTok{values\_to =} \StringTok{"Nb\_Medailles"}
\NormalTok{  )}

\FunctionTok{ggplot}\NormalTok{(data\_graphique, }\FunctionTok{aes}\NormalTok{(}\AttributeTok{x =} \FunctionTok{reorder}\NormalTok{(nationality, }\SpecialCharTok{{-}}\NormalTok{Nb\_Medailles), }\AttributeTok{y =}\NormalTok{ Nb\_Medailles, }\AttributeTok{fill =}\NormalTok{ Type\_Performance)) }\SpecialCharTok{+}
  \FunctionTok{geom\_col}\NormalTok{(}\AttributeTok{position =} \StringTok{"dodge"}\NormalTok{, }\AttributeTok{color =} \StringTok{"white"}\NormalTok{, }\AttributeTok{width =} \FloatTok{0.7}\NormalTok{) }\SpecialCharTok{+}
  \FunctionTok{scale\_fill\_manual}\NormalTok{(}
    \AttributeTok{values =} \FunctionTok{c}\NormalTok{(}\StringTok{"Moyenne\_Historique"} \OtherTok{=} \StringTok{"grey"}\NormalTok{, }\StringTok{"Medailles\_Domicile"} \OtherTok{=} \StringTok{"yellow"}\NormalTok{),}
    \AttributeTok{labels =} \FunctionTok{c}\NormalTok{(}\StringTok{"Moyenne Historique par édition"}\NormalTok{, }\StringTok{"Performance à Domicile"}\NormalTok{)}
\NormalTok{  ) }\SpecialCharTok{+}
  \FunctionTok{labs}\NormalTok{(}
    \AttributeTok{title =} \StringTok{"Comparaison des médailles par édition pour les nations hôtes du dataset"}\NormalTok{,}
    \AttributeTok{x =} \StringTok{"Code Pays (CIO)"}\NormalTok{,}
    \AttributeTok{y =} \StringTok{"Nombre total de médailles par édition"}\NormalTok{,}
    \AttributeTok{fill =} \StringTok{"Indicateur :"}
\NormalTok{  ) }\SpecialCharTok{+}
  \FunctionTok{theme\_minimal}\NormalTok{() }\SpecialCharTok{+}
  \FunctionTok{theme}\NormalTok{(}
    \AttributeTok{plot.title =} \FunctionTok{element\_text}\NormalTok{(}\AttributeTok{face =} \StringTok{"bold"}\NormalTok{, }\AttributeTok{size =} \DecValTok{14}\NormalTok{),}
    \AttributeTok{axis.text.x =} \FunctionTok{element\_text}\NormalTok{(}\AttributeTok{angle =} \DecValTok{45}\NormalTok{, }\AttributeTok{hjust =} \DecValTok{1}\NormalTok{),}
    \AttributeTok{legend.position =} \StringTok{"top"}
\NormalTok{  )}
\end{Highlighting}
\end{Shaded}

\pandocbounded{\includegraphics[keepaspectratio]{rapport_files/figure-latex/unnamed-chunk-16-1.pdf}}

\end{document}
